\chapter{Followers}
\label{ch:followers}
\index{followers}\index{allies}

\begin{tcolorbox}[colback=black!3!white,colframe=black!60!white,title=The Weight of a Name]
The courier taps twice on the shutters of the riverside inn. Inside, your \textbf{Safehouse} (see \ref{ch:compendium-assets}) has already warmed a room, stocked dry clothes, and left an unmarked letter. Outside, your \textbf{Scout} whistles a two-note warning; the watch turns down a side street. Neither scene happens by accident. You paid for it -- in XP, in care, and sometimes in blood.
\end{tcolorbox}

In \textit{Fate's Edge}, your reach is measured by more than sword-arms and spell-forms. \textbf{Assets} move the world when you are elsewhere; \textbf{Followers} act at your side when the camera is on you. Both are power -- and both are promises you will be asked to keep.

This chapter focuses on \textbf{Followers}: the named allies who stand beside you. For Assets (safehouses, charters, networks), see \ref{ch:compendium-assets}.

\section{Follower Theory: Why They Are Expensive}
\label{sec:follower-theory}
\index{followers!theory}\index{costs!followers}

Followers are not hired muscle or disposable hirelings. They are \textbf{named characters with agency, needs, and arcs}. The XP cost (Cap\(^2\)) reflects not merely their skill but the narrative investment required to keep them loyal, safe, and effective.

\paragraph{Why Cap\(^2\) XP?}
\begin{itemize}
  \item \textbf{Cap 1 (1 XP):} A novice or narrow specialist. Cheap because they are replaceable -- but they will not risk their neck without reason.
  \item \textbf{Cap 2 (4 XP):} A competent professional. The cost begins to bite; you will notice if they are hurt.
  \item \textbf{Cap 3 (9 XP):} A skilled veteran. You have shared history; losing them would be a story beat.
  \item \textbf{Cap 4 (16 XP):} An expert with reputation. They are as capable as a low-tier PC.
  \item \textbf{Cap 5 (25 XP):} A master. They could lead their own adventures. You are paying for narrative weight, not just dice.
\end{itemize}

\paragraph{They Act Independently}
Unlike Assets, Followers have \textbf{independent initiative} (once per scene, party-wide). They do not merely add dice; they can change the fiction on their own -- scouting ahead, creating distractions, fetching objects, or shielding you from harm. This action economy is powerful, and the XP cost reflects that.

\paragraph{They Have Needs}
Followers require \textbf{upkeep} (XP or a scene each downtime). Neglect them, and they become \emph{Wary} or \emph{Neglected}; ignore them twice, and they may leave. They have \textbf{personal quests} -- resolve them, and loyalty rises. These are not burdens but story hooks.

\section{Follower Mechanics (Summary)}
\label{sec:follower-mechanics}
\index{followers!mechanics}

\subsection*{Capability Rating (Cap)}
\index{followers!cap}
Followers are rated \textbf{Cap 1--5}. Cost to recruit is \(\textbf{Cap}^2\) XP; 1--3 days of downtime to find, brief, and bind.

\begin{center}
\begin{tabular}{cl}
\toprule
\textbf{Cap} & \textbf{Typical Archetype} \\
\midrule
1 (1 XP) & Novice scout, apprentice, porter \\
2 (4 XP) & Competent guard, clerk, journeyman \\
3 (9 XP) & Veteran mercenary, guild factor, skilled scout \\
4 (16 XP) & Expert spy, master artisan, battle captain \\
5 (25 XP) & Legendary duelist, archivist of forbidden lore, spymaster \\
\bottomrule
\end{tabular}
\end{center}

\textbf{Limit:} You can comfortably maintain up to \(\textbf{Presence} + 1\) Followers without loyalty complications. Exceeding this limit increases the chance of conflict or neglect (GM's discretion).

\subsection*{Recruiting Followers}
\index{followers!recruitment}
Recruitment is a narrative scene, not a menu purchase. The GM should frame a short scene where you encounter a potential follower, negotiate terms, and prove your worth. Common hooks:
\begin{itemize}
  \item You save a life; they offer service to repay a debt.
  \item You share a common enemy; they join to settle a grudge.
  \item They are down on their luck; you offer purpose and pay.
  \item They are an old contact from your background (see \ref{ch:backgrounds}); you call in a favor.
\end{itemize}
Once terms are agreed, you pay XP and downtime as above. The GM may waive the XP cost if the follower's motivation is purely narrative (e.g., a rescued prisoner who becomes a loyal ally). The new follower enters play with \emph{Steady} loyalty and no Exposure or Harm.

\subsection*{Assistance Mechanics}
\index{followers!assistance}
\begin{itemize}
  \item \textbf{Base Assist:} Add dice equal to \(\min(\text{Cap}, \text{relevant Skill})\) to the PC's roll.
  \item \textbf{Hard Cap:} Total assist dice from \emph{all} sources (other PCs, followers, assets) \(\leq\) \textbf{+3d}. \index{dice caps}
  \item \textbf{Specialization:} Some Followers list a \emph{Specialty} granting +1d in a narrow context (still obeys +3d cap).
  \item \textbf{One Helper:} Only one Follower may assist a given action.
\end{itemize}

\subsection*{Condition: Exposure and Harm}
\index{followers!exposure}\index{followers!harm}
Track two meters, each on a 0--4 scale (0=None, 1=Light, 2=Moderate, 3=Severe, 4=Critical). Each level of Moderate or Severe imposes a --1d penalty to the follower's own actions (but not to assist dice). Critical means the follower is out of action (captured, unconscious, fled).
\begin{itemize}
  \item \textbf{Exposure} -- attention, stress, notoriety, overextension.
  \item \textbf{Harm} -- injury, trauma, broken tools, shaken nerve.
\end{itemize}
Followers also have the same \emph{Maintained/Neglected/Compromised} states as Assets (see \ref{ch:compendium-assets}), representing their availability and morale.

\subsection*{Independent Initiative}
\index{followers!initiative}
Once per scene (party-wide), \textbf{one} Follower can act independently instead of assisting. Choose one:
\begin{description}
  \item[Scout \& Signal] Shift an ally's next action to \emph{Dominant} position.
  \item[Distract \& Draw] Reduce a Threat/Hunt/Escape clock by 1 tick.
  \item[Fetch \& Carry] Move an object or person through a hazard safely.
  \item[Cover \& Protect] Eat a consequence meant for a PC; the follower takes Harm 1 or Exposure 2 (GM's choice).
\end{description}
\textbf{Cost:} Mark +1 \emph{Exposure} or \emph{Harm 1} on that Follower as described.

\subsection*{Upkeep and Risks}
\index{followers!upkeep}\index{complications!followers}
\begin{itemize}
  \item \textbf{Maintenance:} Each downtime, pay XP equal to Cap \emph{or} devote a relationship scene. Skip it and they become \emph{Neglected} (lose 1 loyalty step). Two consecutive neglected downturns may cause them to leave.
  \item \textbf{Complication Targeting:} When you gain 2+ SB of consequences on an assisted roll, the GM may assign the fallout to the assisting Follower (fiction first).
  \item \textbf{Off-Screen Actions:} A Cap 3+ follower can attempt a notable off-screen task during downtime; doing so \textbf{generates 1 SB} for the next related scene (the GM determines the nature of the complication).
\end{itemize}

\section{Delegation: Cap 4--5 Followers as Asset Managers}
\label{sec:delegation}
\index{followers!delegation}\index{assets!management}

A Cap 4 or Cap 5 follower is not just a skilled ally -- they are a \textbf{trusted lieutenant} capable of running your off-screen assets while you are away. This frees you from routine maintenance and allows you to focus on adventuring.

\paragraph{The Rule}
If you have a follower with \textbf{Cap 4 or 5}, they may maintain a number of your Assets (see \ref{ch:compendium-assets}) as follows:
\begin{itemize}
  \item A \textbf{Cap 4} follower can maintain up to \textbf{2 Assets}.
  \item A \textbf{Cap 5} follower can maintain up to \textbf{4 Assets}.
\end{itemize}
Maintenance happens without requiring XP or your personal downtime scene. In effect, they \emph{become} the upkeep.

\begin{tcolorbox}[colback=green!5!white,colframe=green!75!black,title=Delegation Example]
You have a Cap 4 Guild Factor follower. You also own four Standard Assets (Guild Seat, Trading Charter, Informant Ring, Safehouse). Normally you would need to spend XP or personal scenes each downtime to keep all four \emph{Maintained}. With delegation, your Cap 4 follower can maintain \textbf{2 of those Assets} for you at no XP cost and without your personal attention. You must maintain the other two yourself (XP or a scene). If you had a Cap 5 follower, they could maintain all four.

If you owned a fifth Asset, you would need to maintain it yourself -- or recruit another Cap 4+ follower.
\end{tcolorbox}

\paragraph{What Delegation Covers}
\begin{itemize}
  \item Routine upkeep (paying staff, restocking supplies, minor repairs).
  \item Handling day‑to‑day communication with faction contacts.
  \item Managing the Asset's condition (preventing it from slipping to \emph{Neglected}).
\end{itemize}

\paragraph{What Delegation Does \textbf{Not} Cover}
\begin{itemize}
  \item \textbf{Crises:} If an Asset is directly threatened (a rival burns your safehouse, your spy ring is compromised), delegation does not automatically resolve the crisis. The follower may alert you, but you must deal with it personally or spend additional resources.
  \item \textbf{Scene Surge:} Delegation does not grant extra \emph{on-screen} activations (still 1 Boon per use). The follower is off‑screen, not in the scene.
  \item \textbf{Multiple Followers:} You cannot combine two Cap 3 followers to maintain a single Asset through delegation -- only Cap 4 or 5 individuals have the authority and expertise to manage holdings independently.
\end{itemize}

\paragraph{GM Guidance}
When a delegated Asset faces a crisis, the GM may:
\begin{itemize}
  \item Advance a \emph{Faction Clock} (see \ref{ch:world-powers-obligation}) representing the threat.
  \item Require the player to spend a downtime scene (or a Boon) to direct the follower's response.
  \item Introduce a new complication: ``Your Guild Factor reports that the warehouse has been seized by customs. What do you do?''
\end{itemize}
Delegation is a privilege, not an immunity. The world still pushes back.

\subsection*{The Risk of Delegation: Rival Factions}
\label{subsec:delegation-risk}
\index{followers!delegation!risks}\index{rival factions}

Delegation is powerful, but it places your Assets in another's hands. A Cap 4--5 follower who becomes \textbf{disgruntled} -- that is, whose loyalty falls to \emph{Wary} or whose condition becomes \emph{Neglected} or \emph{Compromised} -- may decide to keep the Assets they manage.

\paragraph{When Does a Follower Turn?}
The GM may trigger a defection when:
\begin{itemize}
  \item The follower's loyalty drops to \emph{Wary} and you ignore their grievances for two consecutive downtimes.
  \item The follower becomes \emph{Compromised} (captured, blackmailed, or ideologically turned) and the opposition offers a better deal.
  \item You explicitly betray or abandon them (e.g., leave them to take the fall for your crime).
\end{itemize}

\paragraph{What They Take}
A defecting follower does not steal random gear. They take \textbf{the Assets you delegated to them}. They may also recruit other dissatisfied followers or contacts, forming a \textbf{Rival Faction} (see the GM Guide for full rules).

\begin{itemize}
  \item \textbf{Minor Assets (4 XP):} The follower may operate them independently, denying you their use. They do not become enemies overnight -- first they ignore your calls, then they start leveraging the Asset against you (e.g., your safehouse becomes their safehouse).
  \item \textbf{Standard Assets (8 XP):} The follower can convert the Asset into a base of operations for their new faction. You lose access, and the GM starts a \emph{Rival Ambition} [6] clock.
  \item \textbf{Major Assets (12 XP):} The follower may declare themselves the rightful owner (forged deeds, popular support, blackmail). You must contest them in play -- through law, duel, or sabotage.
\end{itemize}

\paragraph{Becoming a Rival Faction}
A defecting Cap 4--5 follower with at least one Asset becomes a \textbf{terrestrial patron} (see \ref{ch:world-powers-obligation}) in their own right. The GM creates a simple faction sheet:
\begin{itemize}
  \item \textbf{Name:} e.g., ``The Sundered Ledger,'' ``The Factor's Faction,'' ``The Broken Chain.''
  \item \textbf{Goal:} What do they want? (Revenge, independence, to prove you wrong.)
  \item \textbf{Assets:} The ones they took (plus any they acquire).
  \item \textbf{Obligation Track:} The party now owes them -- or they owe the party? The GM sets a starting value.
\end{itemize}

The party may:
\begin{itemize}
  \item \textbf{Negotiate:} Pay back wages, offer a public apology, or give up another Asset to buy them off.
  \item \textbf{Intimidate:} Threaten violence or expose secrets to force a return.
  \item \textbf{Take by Force:} Raid the stolen Asset and oust the defector.
  \item \textbf{Let Them Go:} Accept the loss and make a new enemy (which may be more interesting).
\end{itemize}

\paragraph{Preventing Defection}
\begin{itemize}
  \item Keep your Cap 4--5 followers at \emph{Steady} or \emph{Devoted} loyalty.
  \item Address their personal quests (see \ref{sec:follower-personal-quests}) promptly.
  \item Do not leave them in charge of Assets you cannot afford to lose. Spread delegation across multiple loyal followers.
  \item If you sense a follower is unhappy, spend a downtime scene to repair the relationship before they act.
\end{itemize}

\paragraph{Example: The Turned Factor}
Rellan's Cap 4 Guild Factor, Vex, has been \emph{Neglected} for two downturns (Rellan was busy dungeon‑crawling). Vex's loyalty drops to \emph{Wary}. When Rellan returns, Vex presents an ultimatum: ``I have kept your Trading Charter and Informant Ring running, but I am done being ignored. Either give me a seat at the table -- or I take them both.'' Rellan refuses. Vex resigns, takes the two Assets, and founds the ``Free Factors Guild,'' a rival faction that now competes for the same trade routes. Rellan must now decide: crush them, negotiate, or let the enmity simmer.

\begin{tcolorbox}[colback=red!5!white,colframe=red!75!black,title=GM Guidance: Rival Faction from Delegation]
\begin{itemize}
  \item Do not spring defection without warning. Foreshadow with rumours of discontent, missing inventory, or a follower who ``forgets'' to report.
  \item The defecting follower should not become a mustache‑twirling villain unless the table enjoys that tone. They may have legitimate grievances.
  \item The new faction's power should roughly match the party's tier. A Cap 4 follower with one Standard Asset is a credible threat; a Cap 5 with three Major Assets is a campaign villain.
  \item Allow the party to reclaim the Assets through play. Victory feels earned when you storm your own warehouse or argue your case before a magistrate.
\end{itemize}
\end{tcolorbox}

\section{Loyalty and Bonds}
\index{followers!loyalty}\index{bonds!followers}

\subsection*{Loyalty States}
\begin{description}
  \item[Wary] Requires +1 XP upkeep; may refuse dangerous orders.
  \item[Steady] Normal; no extra costs or perks.
  \item[Devoted] Once per arc, downgrade a complication that targets them (Major \(\rightarrow\) Minor). \index{arcs!story}
\end{description}
Loyalty shifts based on treatment, shared victories, and personal quests. A follower who is consistently paid, respected, and protected will rise (Wary → Steady after 2--3 sessions; Steady → Devoted after completing a personal quest). One who is ignored, underpaid, or sacrificed will fall.

\subsection*{Bond Integration}
\index{bonds}\index{boons!bonds}
If a Follower is also a \textbf{Bond} (see \ref{ch:backgrounds}), they grant the bonded Boon (once per session, +1 Boon when you act to uphold the bond) in addition to their assist dice. Working together to resolve the follower's personal quest counts as upholding the bond.

\subsection*{Follower Personal Quests}
\label{sec:follower-personal-quests}
\index{followers!personal quests}
Every follower has a hidden (or open) goal. Discover it through play. When you help them achieve it, they:
\begin{itemize}
  \item Gain a permanent loyalty boost (Steady → Devoted).
  \item May unlock a new specialization (GM's discretion).
  \item You earn a \textbf{Legacy Point} (see \ref{ch:legacy-engine} for use).
\end{itemize}
Examples:
\begin{itemize}
  \item Avenge a murdered family member.
  \item Recover a lost heirloom or debt note.
  \item Prove themselves worthy of a guild title.
  \item Find a cure for a chronic illness.
  \item Escape a past that still hunts them.
\end{itemize}

\section{Advanced Follower Management}
\index{followers!advancement}\index{followers!specialization}

\subsection*{Advancement}
On significant shared accomplishments, a Follower may increase Cap by 1 (max PC tier -- 1). XP cost equals \(\text{new Cap}^2 - \text{current Cap}^2\). One advancement per major arc.

\subsection*{Common Specializations}
\begin{description}
  \item[Combat Specialist] +1d assisting violent actions; may duel Cap\(\leq\)3 foes.
  \item[Social Specialist] +1d on parley, rumor-work, courtly rites.
  \item[Technical Specialist] +1d on craft/repair; may clear \emph{Compromised} on one item per downtime.
  \item[Magic Specialist] +1d assisting \emph{Weave} (see \ref{ch:magic}); may absorb 1 SB of backlash once/session. \index{magic!casting loop}
\end{description}

\subsection*{Groups and Cohorts}
\index{followers!groups}
A squad, crew, or circle can be run as a single \emph{Group Follower} with one Cap, one Exposure track (shared), and distributed Harm. Treat as one assistant in play. When the group takes Harm, it represents casualties or loss of cohesion; the group may be reduced in Cap temporarily or permanently.

\subsection*{Retirement and Legacy}
\index{followers!retirement}\index{legacy engine}
A follower may retire from active duty after completing their personal quest or suffering critical harm. They become an \textbf{Off-Screen Asset} (Minor for Cap 1--2, Standard for Cap 3+, free of XP cost) representing their ongoing network or expertise. They may also become a new PC via the Legacy Engine (see \ref{ch:legacy-engine}). See the \textbf{Legacy Planning Checklist} (see \ref{sec:legacy-checklist}) for steps to prepare a follower as a successor.

\section{Sample Followers}
\index{followers!examples}

\subsection*{Cap 1: Street Informant}
\textbf{Skills:} Stealth 2, Sway 1, Notice 1 \quad
\textbf{Assist:} +1d to urban intel \\
\textbf{Specialty:} Knows back ways; reduces escape clock DV by 1 once/session. \\
\textbf{Quirk:} Paid in favors and gossip; hates written notes. \\
\emph{Vignette:} \textit{``Two taps on the shutters if it's Watch. Three if it's worse.''}

\subsection*{Cap 2: Trusted Bodyguard}
\textbf{Skills:} Melee 2, Athletics 2, Notice 2 \quad
\textbf{Assist:} +2d on protection/combat \\
\textbf{Specialty:} Can \emph{Cover \& Protect} without taking Harm once/session. \\
\textbf{Quirk:} Never speaks first; always positions themselves between you and exits.

\subsection*{Cap 3: Veteran Mercenary}
\textbf{Skills:} Melee 3, Tactics 2, Athletics 2 \quad
\textbf{Assist:} +3d on combat/tactics \\
\textbf{Specialty:} Holds lines; can \emph{Cover \& Protect} without taking Harm once/session. \\
\textbf{Quirk:} Won't fight for ``poetry'' or ``exposure.''

\subsection*{Cap 4: Guild Factor}
\textbf{Skills:} Sway 3, Investigation 3, Lore 2 \quad
\textbf{Assist:} +3d on commerce/legal matters \\
\textbf{Specialty (Delegation):} Can maintain up to 2 Assets without PC upkeep (see \ref{sec:delegation}). \\
\textbf{Quirk:} Keeps meticulous records; remembers every favor and debt.

\subsection*{Cap 5: Master Artificer}
\textbf{Skills:} Craft 4, Lore 3, Investigation 3 \quad
\textbf{Assist:} +3d on crafting/research \\
\textbf{Specialty (Delegation):} Can maintain up to 4 Assets without PC upkeep (see \ref{sec:delegation}). \\
\textbf{Quirk:} Collects flawed prototypes ``for teaching.''

\section{Strategic Play with Followers}

\subsection*{Choosing Followers}
\begin{itemize}
  \item \textbf{Diversity beats duplication:} one combat, one social, one technical.
  \item \textbf{Reliability over raw Cap:} A steady Cap 2 is worth more than a wary Cap 4.
  \item \textbf{Spread risk:} Rotate who takes Exposure to avoid Critical spikes.
  \item \textbf{Bond synergy:} Recruit followers who share your character's Bonds for bonus Boon generation.
\end{itemize}

\subsection*{Keeping Followers Alive}
\begin{itemize}
  \item \textbf{Use Cover \& Protect sparingly:} It keeps you alive but burns followers.
  \item \textbf{Clear Exposure early:} A single downtime scene can reduce Exposure from 2 to 0.
  \item \textbf{Invest in healing:} A follower with Harm 2 is a liability. Pay for treatment or use a magic specialist.
  \item \textbf{Know when to retreat:} Losing a follower is worse than losing a battle.
\end{itemize}

\subsection*{Leveraging Independent Initiative}
\begin{itemize}
  \item \textbf{Scout \& Signal is free Position:} Use it before a big fight or chase.
  \item \textbf{Distract \& Draw on clocks:} Reduce a pursuit or alarm clock by 1 -- often worth more than a +1d.
  \item \textbf{Fetch \& Carry can solve puzzles:} Move a key item across a chasm while you fight.
  \item \textbf{Cover \& Protect as last resort:} Better than you dying, but expensive.
\end{itemize}

\section{Asset \& Follower Synergy}
\index{synergy!assets and followers}
\begin{itemize}
  \item \textbf{Workshop + Artificer:} Clear \emph{Compromised} on gear automatically once/session.
  \item \textbf{Spy Ring + Informant:} First recon opener starts \emph{Dominant} \emph{and} yields a concrete lead.
  \item \textbf{Temple Chapter + Social Specialist:} Convert a \emph{Crowd} complication into a \emph{Petition} clock.
  \item \textbf{Healing House + Combat Specialist:} Recover 1 Harm on a follower between sessions.
\end{itemize}

\section{Integration with Core Systems}
\index{SB economy}\index{boons}\index{downtime}
\begin{itemize}
  \item \textbf{SB Economy:} Asset use under pressure invites SB spends; Followers can catch those consequences if fictionally exposed.
  \item \textbf{Boons:} Spend to activate an Asset in-scene, rally a Follower from \emph{Moderate} to \emph{Light} condition, or fuel an Independent Initiative.
  \item \textbf{Downtime:} Best place to repair, recruit, advance, and negotiate terms. Tie these scenes to Bonds to harvest story.
\end{itemize}

\section{GM Tools for Followers}

\subsection*{Quick SB Menu (Followers)}
\begin{itemize}
  \item \textbf{1 SB:} Mark \emph{Exposure 1} -- they are noticed by the wrong eyes.
  \item \textbf{2 SB:} \emph{Harm 1} or \emph{Exposure 2} -- their nerves fray or they are bruised.
  \item \textbf{3 SB:} Lose a tool or contact; --1d until replaced.
  \item \textbf{4 SB:} \emph{Compromised} -- captured, spooked, or defecting under pressure (but not killed unless the table has established that as a risk).
\end{itemize}

\subsection*{Follower Personal Quest Seeds}
\index{adventure seeds!follower quests}
\begin{itemize}
  \item \textbf{The Lost Ledger:} Your factor's father disappeared with a book of debts. Find it before his rivals do.
  \item \textbf{The Broken Oath:} Your mercenary deserted a previous contract; now a bounty hunter seeks them. Protect them or negotiate a pardon.
  \item \textbf{The Silent Bell:} Your scout's family village stopped answering the three-tone law. Something in the mist took them.
  \item \textbf{The Artificer's Rival:} The master artificer's old apprentice is selling flawed enchantments under their name. Clear the reputation.
\end{itemize}

\subsection*{Challenging Followers Without Killing Them}
\begin{itemize}
  \item \textbf{Capture instead of kill:} Rivals want leverage, not corpses.
  \item \textbf{Conditional harm:} ``Your arm is broken; you cannot assist in combat until healed.''
  \item \textbf{Moral dilemmas:} The follower is offered a bribe to betray you; they resist but are shaken (Exposure +2).
  \item \textbf{Leave clues:} A missing follower leaves a trail the party can follow.
\end{itemize}

\section{Narrative Seeds}
\index{adventure seeds}\index{hooks!assets}\index{hooks!followers}
\paragraph{The Broken Ledger.}
Your \emph{Banking Charter} flags a forged letter-of-mark. Do you protect the account or the client's life?

\paragraph{The Bell at Dusk.}
A Mistlands \emph{Safehouse} stops answering the three-tone law. The bell rope has been cut from the outside.

\paragraph{The White Flag.}
Your Cap 3 mercenary refuses to take a fortified bridge without pay in hand. He is right -- someone is playing both sides.

\paragraph{The Informant's Debt.}
Your street informant calls in a favor from years ago. They need you to intimidate a landlord -- or someone will die.

\section*{Quick Reference}
\begin{tcolorbox}[colback=blue!5!white,colframe=blue!75!black,title=Followers Quick Reference,fonttitle=\bfseries]
\textbf{Follower Basics}
\begin{itemize}
  \item Cost: Cap\(^2\) XP; Limit: \textbf{Presence} + 1 without strain.
  \item Assistance: +\(\min(\)Cap, Skill\()\) dice, \textit{max +3d total from all sources}.
  \item Independent Action (party 1/scene): Scout \& Signal; Distract \& Draw; Fetch \& Carry; Cover \& Protect (cost: +Exposure or Harm 1).
  \item Tracks: Exposure/Harm 0--4; states mirror Assets (Maintained/Neglected/Compromised).
  \item Loyalty: Wary (costs +1 XP upkeep), Steady, Devoted (once/arc downgrade complication).
\end{itemize}

\textbf{Delegation (Cap 4--5 only)}
\begin{itemize}
  \item A Cap 4 follower can maintain up to \textbf{2 Assets} without XP or personal downtime scenes.
  \item A Cap 5 follower can maintain up to \textbf{4 Assets}.
  \item Does not cover crises (rival attacks, theft) -- those still require your attention.
  \item Cannot be combined from multiple lower‑Cap followers.
\end{itemize}
\end{tcolorbox}

\section*{Closing Image}
\index{flourishes}
\textit{A shuttered lamp, a sealed ledger, a scarred gauntlet resting on a map. You do not move pieces -- you invest in people and places. When you step into the street, the city already leans your way.}